\documentclass[12pt]{article}
\usepackage[utf8]{inputenc}
\usepackage[T1]{fontenc}
\usepackage{hyperref}
\title{Ventureño Language Resources for K-12 Curriculum Development}
\author{Compiled by Claude Sonnet 5 (claude-sonnet-5) for Andr\'as Kornai}
\date{}
\begin{document}
\maketitle

\section{Introduction}
Ventureño (ISO 639-3 veo; Chumashan) is dormant. Sources differ on its status. One lists it as extinct, while a Chumash community dictionary site says the Mitsqanaqan' dialect went to sleep in the early 1960s and describes current revitalization efforts. Search turned up no school curriculum or textbook. The material is archival and lexical, and it is mostly built on John P. Harrington's fieldnotes.

\section{Primary Educational Resources}
No primers, textbooks or classroom materials turned up in search. A curriculum would have to be built from the dictionary below.

\section{Linguistic and Reference Materials}
\begin{itemize}
\item Timothy Henry-Rodriguez, \emph{A Mitsqanaqan' Ventureño--English, English--Mitsqanaqan' Ventureño Dictionary} (2019). It has about 5000 Ventureño entries, an English-to-Ventureño section, and semantic word lists. Entries draw on speakers and on the Harrington archives. Open-access record: \url{https://scholarspace.manoa.hawaii.edu/collections/bf4ed241-69d9-477a-81a3-1605d499316f}
\item The same dictionary on the UBC wiki: \url{https://wiki.ubc.ca/Documentation:RelLex/A_Mitsqanaqan%CC%93_Venture%C3%B1o_%E2%80%94_English,_English_%E2%80%94_Mitsqanaqan%CC%93_Venture%C3%B1o_Dictionary}
\item CSU Channel Islands Chumash dictionary site, with a language-family page: \url{https://ciapps.csuci.edu/ChumashDictionary/About}
\end{itemize}

\section{Religious and Community Resources}
None found in search.

\section{Digital and Audio Resources}
The Smithsonian National Anthropological Archives hold Harrington materials, including an English--Chumash vocabulary manuscript (MS 858) and Chumash sound recordings. Finding aid: \url{https://sirismm.si.edu/EADpdfs/NAA.MS858.pdf}

\section{Curriculum Development Context}
Little exists beyond the dictionary and the archives. Any K-12 programme would be a revitalization effort from lexical sources, and it should work with the Chumash community.
\end{document}